\metatitle{Posets and linear extensions}
\metadescription{Finite posets, linear extensions, Jordan--Hölder sets, order ideals, hook formulas, order polytopes, and common classes of posets.}
\metakeywords{posets,linear extensions,Jordan-Hölder set,order ideals,order polytopes,series-parallel posets,lattices,3+1-free posets}

\section[posets]{Posets}

A \defin{poset}\label{poset} is a set $P$ with a partial order $\leq_P$:
the relation is reflexive, antisymmetric, and transitive.  We write
$x\lt_P y$ if $x\leq_P y$ and $x\neq y$.  A
\defin{cover relation}\label{coverRelation} is a relation
$x\lessdot_P y$ such that $x\lt_P y$ and there is no $z$ with
$x\lt_P z\lt_P y$.  The
\defin{Hasse diagram}\label{hasseDiagram} of $P$ draws only these cover
relations.

Two basic examples are the \defin{chain}\label{chainPoset}
$1\lt 2\lt \dotsb \lt n$ and the \defin{antichain}\label{antichain} on $n$
elements, where no two distinct elements are comparable.  Posets are useful
because many combinatorial objects are naturally sets with precedence
constraints: boxes in a Young diagram, roots in a root poset, subsets ordered
by inclusion, or independent sets ordered by inclusion.

\section[linearExtensions]{Linear extensions}

Let $P$ be a finite poset with $n$ elements.  A
\defin{linear extension}\label{linearExtension} of $P$ is a total ordering
$(p_1,p_2,\dotsc,p_n)$ of the elements of $P$ such that
\[
  p_i\lt_P p_j \quad\Longrightarrow\quad i\lt j.
\]
Equivalently, it is an order-preserving bijection $P\to [n]$.
We write $e(P)$ for the number of linear extensions of $P$.

\begin{example*}[Three small posets]
Let $P$ be the poset with two minimal elements $a,b$ and one maximal element
$c$, so that $a\lt c$ and $b\lt c$.  Its linear extensions are
\[
  (a,b,c),\qquad (b,a,c).
\]
For the chain $a\lt b\lt c$, there is exactly one linear extension:
$(a,b,c)$.  For the antichain on $\{a,b,c\}$, every ordering is allowed, so
there are $3!=6$ linear extensions.
\end{example*}

Counting linear extensions is difficult in general: \name{Graham Brightwell}
and \name{Peter Winkler} prove that this counting problem is
\#P-complete \cite{BrightwellWinkler1991}.  This is one reason why product
formulas, recurrences, and special poset classes are valuable.

\section[jordanHolderSet]{Jordan--Hölder sets}

A linear extension is usually a word in the \emph{elements} of $P$.  A
\defin{Jordan--Hölder set}\label{jordanHolderSet} records instead the
\emph{labels} that appear along a linear extension.

More precisely, let $w:P\to [n]$ be a bijective labeling.  The
\defin{labeled poset}\label{labeledPoset} is the pair $(P,w)$.  A labeling is
\defin{natural}\label{naturalLabeling} if $x\lt_P y$ implies
$w(x)\lt w(y)$.  The Jordan--Hölder set of the labeled poset $(P,w)$ is
\[
  \mathcal{L}(P,w)
  \coloneqq
  \{\sigma\in\symS_n : \sigma^{-1}\circ w \text{ is order-preserving}\}.
\]
Thus $\sigma$ is the word of labels seen in some linear extension of $P$.
The poset determines the possible element orders; the labeling translates
those orders into permutations.

\begin{example*}[Same poset, different labels]
Let $P$ be the poset $a,b\lt c$.  With the natural labeling
$w(a)=1$, $w(b)=2$, and $w(c)=3$, the Jordan--Hölder set is
\[
  \mathcal{L}(P,w)=\{123,213\}.
\]
If we relabel by $w(c)=1$, $w(a)=2$, and $w(b)=3$, then the linear extensions
of the underlying poset are still $(a,b,c)$ and $(b,a,c)$, but the
Jordan--Hölder set becomes
\[
  \mathcal{L}(P,w)=\{231,321\}.
\]
This distinction is important in the theory of
\hyperref[pPartition]{$P$-partitions}, where descent sets of the label words
control the expansion in the \hyperref[gessel]{fundamental quasisymmetric}
basis.
\end{example*}

\section[pEulerianPolynomial]{$P$-Eulerian polynomials}

Let $(P,w)$ be a labeled poset.  The
\defin{$P$-Eulerian polynomial}\label{pEulerianPolynomial} is the descent
generating polynomial of the Jordan--Hölder set:
\[
  W_{P,w}(t)
  \coloneqq
  \sum_{\sigma\in\mathcal{L}(P,w)} t^{|\DES(\sigma)|}.
\]
If the labeling is understood, or if $P$ is naturally labeled, we often write
$W_P(t)$.

\begin{example*}
For the naturally labeled poset $1,2\lt 3$, the Jordan--Hölder set is
$\{123,213\}$.  Thus
\[
  W_P(t)=1+t,
\]
because $123$ has no descent and $213$ has one descent.  With the relabeling
from the previous example, the Jordan--Hölder set is $\{231,321\}$, so the
corresponding labeled-poset Eulerian polynomial is $t+t^2$.
\end{example*}

The $P$-Eulerian polynomial appears in the Ehrhart theory of
\hyperref[orderPolytope]{order polytopes}: the $h^*$-polynomial of
$\mathcal{O}(P)$ is the $P$-Eulerian polynomial for a natural labeling.  This
is the bridge between \hyperref[linearExtension]{linear extensions},
\hyperref[descentSet]{descent statistics}, and \hyperref[posetRealRooted]{real-rootedness
questions for posets}.  A famous question in this direction was the
\hyperref[neggersStanleyConjecture]{Neggers--Stanley conjecture}, which asked
for real-rootedness of these polynomials in broad generality.

\todo{Add a small table comparing $W_{P,w}(t)$ for two non-isomorphic
three- or four-element labeled posets.}

\section[posetIdeals]{Order ideals, filters and antichains}

An \defin{order ideal}\label{orderIdeal} of a poset $P$ is a subset
$I\subseteq P$ such that $x\in I$ and $y\leq_P x$ imply $y\in I$.  Dually, an
\defin{order filter}\label{orderFilter} is an upward-closed subset.  An
\defin{antichain} is a subset whose elements are pairwise incomparable.

Order ideals and antichains are equivalent data: an ideal is determined by
its maximal elements, and those maximal elements form an antichain.  This is
the starting point for many actions such as rowmotion and promotion, and it is
also why the lattice of ideals of a poset appears throughout algebraic
combinatorics.

\begin{example*}[Ideals in a three-element poset]
For the poset $a,b\lt c$, the order ideals are
\[
  \emptyset,\qquad \{a\},\qquad \{b\},\qquad \{a,b\},\qquad \{a,b,c\}.
\]
The maximal elements of these ideals are respectively
\[
  \emptyset,\qquad \{a\},\qquad \{b\},\qquad \{a,b\},\qquad \{c\},
\]
which are precisely the antichains of the poset.
\end{example*}

Independent sets of a graph or matroid can also be ordered by inclusion.  For
a graph $G$, the independent sets form a down-set in the Boolean lattice
$2^{V(G)}$.  This poset is usually not a lattice under union, but it is a
natural way to turn graph-theoretic independence into an ordered object.

\section[posetHookFormulas]{Hook formulas}

Many important posets have unexpectedly simple formulas for $e(P)$.

For a partition $\lambda$, the boxes of its Young diagram form a poset by
declaring a box to be less than the boxes weakly southeast of it.  A standard
Young tableau of shape $\lambda$ is the same thing as a linear extension of
this Young-diagram poset.  The classical
\hyperref[partitionCores]{hook formula} says
\[
  e(\lambda)=|\SYT(\lambda)|
  = \frac{|\lambda|!}{\prod_{\square\in\lambda} h(\square)}.
\]
Equivalently, this is the Frame--Robinson--Thrall hook-length formula
\cite{FrameRobinsonThrall}.  See also the
\hyperref[prelimQHook]{$q$-hook formula}.

There is an analogous hook formula for rooted trees.  If $T$ is a rooted tree
viewed as a poset with $u\leq_T v$ when $u$ is on the path from the root to
$v$, then
\[
  e(T)=\frac{|T|!}{\prod_{v\in T} |T_v|},
\]
where $T_v$ is the rooted subtree consisting of $v$ and all its descendants.
This formula is discussed, for example, by \name{Donald E. Knuth}
\cite{Knuth1998ArtOfProgramming}.

Hook formulas also exist for shifted shapes, $d$-complete posets, and several
other structured families.  The important contrast is that such formulas are
special: for an arbitrary poset, counting linear extensions is computationally
hard \cite{BrightwellWinkler1991}.

\section[posetPolytopes]{Order and chain polytopes}

Finite posets have a direct polyhedral avatar.  The
\hyperref[orderPolytope]{order polytope} $\mathcal{O}(P)$ is the set of
functions $f:P\to [0,1]$ such that $x\leq_P y$ implies $f(x)\leq f(y)$.
The \hyperref[chainPolytope]{chain polytope} $\mathcal{C}(P)$ is cut out by
nonnegativity and the inequalities
\[
  f(x_1)+f(x_2)+\dotsb+f(x_k)\leq 1
\]
for every chain $x_1<_P x_2<_P\dotsb<_P x_k$.  Stanley introduced these two
poset polytopes in \cite{Stanley86TwoPosetPolytopes}.

The normalized volume of $\mathcal{O}(P)$ is $e(P)$, and the
\hyperref[ehrhart]{$h^*$-polynomial} of $\mathcal{O}(P)$ is the descent
generating polynomial of the linear extensions of $P$.  Thus order polytopes
connect linear extensions to Ehrhart theory and to real-rootedness questions
for descent polynomials.

The same story appears in the theory of
\hyperref[pPartition]{$P$-partitions}: the
\hyperref[orderPolynomials]{order polynomial} $\Omega_P(m)$ counts
order-preserving maps $P\to [m]$, and its generating function has numerator
the $P$-Eulerian polynomial.

\section[posetDynamics]{Promotion and rowmotion}

Posets also carry natural dynamics.  \hyperref[promotion]{Promotion} can be
defined on linear extensions of a poset; see \name{Richard Stanley}
\cite{Stanley2009}.  For arbitrary labelings, \name{Colin Defant} and
\name{Noah Kravitz} introduce promotion sorting \cite{DefantKravitz2022}.
Promotion is also one of the standard cyclic actions behind
\hyperref[cyclicSieving]{CSPs} for \hyperref[cspRectangularSYT]{rectangular
standard Young tableaux} and \hyperref[cspKrewerasWords]{Kreweras words}.

\begin{example*}[Promotion on a small poset]
Let $P$ be the poset $a,b\lt c$.  The two linear extensions are
\[
  (a,b,c),\qquad (b,a,c).
\]
Under promotion, the first extension follows the chain $a\lt c$ and becomes
$(b,a,c)$.  Similarly, $(b,a,c)$ follows the chain $b\lt c$ and becomes
$(a,b,c)$.  Thus promotion has one orbit of size $2$ on this example.
\end{example*}

On order ideals, one important action is \defin{rowmotion}\label{rowmotion}:
send an ideal to the ideal generated by the minimal elements outside it.  Root
posets and minuscule posets give many cyclic-sieving examples; see the
\hyperref[cspListPosetsRoot]{root poset} and
\hyperref[cspListPosetsMinuscule]{minuscule poset} sections on cyclic
sieving.  Rowmotion and its toggle variants are also central examples in
\hyperref[homomesy]{homomesy}.

\begin{example*}[Rowmotion on a chain]
Let $P$ be the chain $x\lt y\lt z$.  Its order ideals form one rowmotion
orbit:
\[
  \emptyset
  \longmapsto \{x\}
  \longmapsto \{x,y\}
  \longmapsto \{x,y,z\}
  \longmapsto \emptyset.
\]
Indeed, from $\{x\}$ the minimal element outside the ideal is $y$, and the
ideal generated by $y$ is $\{x,y\}$.
\end{example*}

\section[posetClasses]{Some classes of posets}

The \defin{width}\label{posetWidth} of a poset is the largest size of an
antichain.  Width-one posets are chains; width-two posets already contain
rich examples such as the posets attached to skew Ferrers diagrams and
\hyperref[latticePathMatroids]{lattice path matroids}.

A \defin{lattice}\label{latticePoset} is a poset in which every two elements
have a meet and a join, usually written $x\wedge y$ and $x\vee y$.  A basic
source of lattices is the set $J(P)$ of order ideals of a finite poset $P$,
ordered by inclusion.  These are distributive lattices.

A \defin{series-parallel poset}\label{seriesParallelPoset} is built from
one-element posets using two operations: disjoint union, also called parallel
composition, and ordinal sum, also called series composition.  Such posets are
often amenable to recursive enumeration because both operations have simple
effects on linear extensions.  These operations preserve real-rootedness of
$P$-Eulerian polynomials; see the \hyperref[posetRealRooted]{real-rootedness
discussion}.

A poset is \defin{$(3+1)$-free}\label{threePlusOneFreePoset} if it has no
induced subposet consisting of a three-element chain together with an element
incomparable to all three elements of the chain.  These posets are central in
the study of incomparability graphs and chromatic symmetric functions:
\name{Vesselin Gasharov} proved Schur positivity for incomparability graphs
of $(3+1)$-free posets \cite{Gasharov1996}.  They also occur in
\hyperref[chromaticQuasisymmetric]{chromatic quasisymmetric functions}.

A poset is \defin{$(2+2)$-free}\label{twoPlusTwoFreePoset} if it has no
induced subposet consisting of two disjoint two-element chains.  Equivalently,
it is an \defin{interval order}\label{intervalOrder}: its elements can be
represented by intervals on the real line, ordered by lying completely to the
left.

A \defin{unit interval order}\label{unitIntervalOrder} is an interval order
that can be represented using intervals of equal length.  Equivalently, it is
both $(3+1)$-free and $(2+2)$-free.  Its incomparability graph is a
\hyperref[chromaticQuasisymmetricUnitIntervalGraph]{unit interval graph}.
This is the setting of many
\hyperref[chromaticQuasisymmetricSchurExpansion]{Schur-positivity} results for
chromatic quasisymmetric functions, and the reduced setting for the
\hyperref[chromaticEExpansionConjectures]{Stanley--Stembridge
$\elementaryE$-positivity problem}.

\todo{Add one explicit $(3+1)$-free example and one induced $(3+1)$
obstruction, preferably with a small Hasse diagram.}
